\section{Background and Problem Grounding}
\label{sec:background}

\subsection{Cloud-Native Telemetry Modalities: Metrics, Logs, and Traces}
In distributed architectures orchestrated by Kubernetes, observability relies on three complementary data modalities:
\begin{enumerate}[leftmargin=1.2em, itemsep=2pt, topsep=3pt]
    \item \textbf{Numerical Metrics ($\mathcal{M}$):} Numeric, timestamped aggregations of operational indicators (e.g., container CPU millicores, cgroup memory working sets, network socket drop rates). Metrics are computationally efficient to scrape, store, and query via domain-specific engines like Prometheus (TSDB) using PromQL. However, metrics lack individual transaction context and request-level execution causality.
    \item \textbf{Structured Logs ($\mathcal{L}$):} Semi-structured, discrete text events generated by application runtimes and container engines (e.g., standard output, stack traces, \texttt{SIGTERM} shutdown notices). Modern log aggregators such as Grafana Loki index streams by metadata labels (namespace, pod name) rather than full-text indexing, optimizing storage footprint while allowing regex-based log streaming via LogQL.
    \item \textbf{Distributed Traces ($\mathcal{T}$):} Directed Acyclic Graphs (DAGs) representing end-to-end execution paths of distributed requests traversing multi-tier microservices \cite{sigelman_2010}. Propagated across network hops via standardized metadata headers (e.g., W3C TraceContext), traces consist of individual spans denoting timed units of work with parent-child relationships.
\end{enumerate}

\subsection{The Diagnostic Dilemma in Unimodal Observability}
Although modern microservices instrument all three pillars, existing root-cause analysis (RCA) techniques typically process them in isolation, triggering three fundamental pathologies:

\subsubsection{The Cascading Span Bias in Distributed Tracing}
Consider an upstream edge gateway ($S_{\text{edge}}$) invoking a chain of internal services $S_1 \rightarrow S_2 \rightarrow S_3$, where $S_3$ experiences a severe network partition or $2000\,\text{ms}$ socket timeout. Because parent spans synchronously enclose the latency of all downstream child spans:
\begin{equation}
\begin{split}
    T_{\text{inclusive}}(S_{\text{edge}}) &> T_{\text{inclusive}}(S_1) \\
    &> T_{\text{inclusive}}(S_2) \ge T_{\text{inclusive}}(S_3)
\end{split}
\end{equation}
Any naive diagnostic algorithm that evaluates span degradation solely by absolute span duration will flag $S_{\text{edge}}$ as the slowest service in the trace. We formally designate this pathology as the \textbf{Cascading Span Bias}. In our baseline experiments, this bias induces a catastrophic $0.00\%$ root-cause accuracy, as the edge gateway is consistently and erroneously indicted for downstream faults.

\subsubsection{The Silent Metric Trap in Infrastructure Monitoring}
Infrastructure metrics are scraped asynchronously at coarse intervals (typically $15\,\text{s}$ to $60\,\text{s}$ to mitigate scraper overhead). When a container experiences a short-lived micro-burst or non-fatal thread contention, the aggregated metric averages out the perturbation. Conversely, when a container terminates abruptly due to a \texttt{SIGKILL} signal, metric scrapers simply lose the target, resulting in missing data rather than an explicit anomaly signal. Without correlated container crash logs, metrics alone cannot pinpoint the failure trigger.

\subsubsection{The Observability Tax and Invasiveness Envelope}
Injecting deep in-process instrumentation hooks across polyglot microservice runtimes (Java Virtual Machines, Node.js V8 engines, Python interpreters, Go binaries) introduces significant CPU and memory bloat, thread contention, and garbage collection pauses \cite{vancamp24, valbusa_2025}. An unconstrained observability pipeline can consume over $30\%$ of total node memory in resource-constrained edge or cloud deployments, transforming the monitoring infrastructure itself into a primary source of cluster instability.

\subsection{Guiding Control Principles for Multimodal Observability}
To overcome these limitations, we formulate two fundamental engineering control principles that govern the PANOPTES reference architecture:
\begin{itemize}[leftmargin=1.2em, itemsep=2pt, topsep=3pt]
    \item \textbf{Control Principle 1 (CP1 — Least Invasive Observation):} Telemetry collection must maximize diagnostic entropy reduction while strictly bounding the dedicated observability control-plane footprint within a predictable resource envelope ($\le 10\%$ CPU, $\le 20\%$ RAM of the monitored workload).
    \item \textbf{Control Principle 2 (CP2 — Semantic Continuity):} Causal context must be deterministically preserved across asynchronous IPC and protocol boundaries by decoupling parent-child inclusive latencies and binding discrete logs and metric time series directly to trace spans via OpenMetrics Exemplars.
\end{itemize}
